Le prime misure effettuate in laboratorio sono state le varie distanze che compie il fascio di luce laser. 
Le lunghezze delle misure $\overline{M_rS_1}$, $\overline{S_1S_2}$ e $\overline{S_1M_f}$ sono riportate nella tabella sottostante (\ref{parziali di D}):
 \begin{table}[ht]
    \centering
    \begin{tabular}{|c|c|c|}
    \hline
     $\overline{M_rS_1}$ [m] & $\overline{S_1S_2}$ [m] & $\overline{S_1M_f}$ [m]\\
    \hline
    $6.06 \pm\, 0.01$ & $0.66 \pm\, 0.01$ & $6.60 \pm\, 0.01$ \\
    \hline
    \end{tabular}
    \caption{In questa tabella sono riportate le lunghezze delle varie componenti che costituiscono il tragitto del fascio dallo specchio rotante allo specchio concavo }
    \label{parziali di D}
\end{table}




Nella tabella sottostante (\ref{primidati}) la grandezza $f_2$ rappresenta la distanza focale della lente $L_2$, la grandezza $D$ rappresenta la distanza tra lo specchio rotante e lo specchio concavo, la grandezza $a$ rappresenta la distanza tra la lente $L_2$ e lo specchio rotante.
\begin{table}[ht]
    \centering
    \begin{tabular}{|c|c|c|}
    \hline
     $f_2$[m] & $D$ [m] & $a$[m] \\
    \hline
    $0.252\pm0.001$ & $13.32\pm0.02$ & $0.451\pm0.001$ \\
    \hline
    \end{tabular}
    \caption{In questa tabella sono riportate le misure preliminari effettuate in laboratorio sull'apparato con le relative incertezze: $f_2$ è la lunghezza focale della lente $L_2$, $D$ è la distanza tra lo specchio rotante e lo specchio concavo e $a$ è la distanza tra la lente $L_2$ e lo specchio rotante.}
    \label{primidati}
\end{table}
\subsection{Senso antiorario}
Il primo set di misure prese sono state effettuate facendo girare lo specchio in senso antiorario:
\begin{table}[ht]
    \centering
    \begin{tabular}{|c|c|c|c|c|}
    \hline
     $\nu_0$ [$Hz$] & $\nu$ [$Hz$] & $\Delta \omega$ [1/s] & $\Delta\delta[m]$ & $c_{spe}$ [$m/s$] \\
    \hline
    $105$ & $282$ & $1112.12$ & $4\times10^{-5}$ & $3.678\times10^8$ \\
    \hline
    $643$ & $860$ & $1363.45$ & $7\times10^{-5}$ & $2.577\times10^8$ \\
    \hline
    $271$ & $426$ & $973.89$ & $4\times10^{-5}$ & $3.221\times10^8$ \\
    \hline
    $123$ & $522$ & $2506.99$ & $1.2\times10^{-4}$ & $2.764\times10^8$ \\
    \hline
    $522$ & $906$ & $2412.74$ & $1.2\times10^{-4}$ & $2.659\times10^8$ \\
    \hline
    $906$ & $1004$ & $615.75$ & $3\times10^{-5}$ & $2.715\times10^8$ \\
    \hline
    $69$ & $336$ & $1677.611$ & $7\times10^{-5}$ & $3.171\times10^8$ \\
    \hline
    $336$ & $611$ & $1727.88$ & $7\times10^{-5}$ & $3.265\times10^8$ \\
    \hline
    $611$ & $886$ & $1727.87$ & $8\times10^{-5}$ & $2.857\times10^8$ \\
    \hline
    $1004$ & $778$ & $1419.99$ & $6\times10^{-5}$ & $3.131\times10^8$ \\
    \hline
    $778$ & $482$ & $1859.82$ & $8\times10^{-5}$ & $3.075\times10^8$ \\
    \hline
    $1014$ & $735$ & $1753.01$ & $7\times10^{-5}$ & $3.313\times10^8$ \\
    \hline
    $735$ & $432$ & $1903.81$ & $1\times10^{-4}$ & $2.519\times10^8$ \\
    \hline
    $150$ & $897$ & $4693.54$ & $2.1\times10^{-4}$ & $2.957\times10^8$ \\
    \hline
    $817$ & $353$ & $2915.39$ & $1.30\times10^{-4}$ & $2.967\times10^8$ \\
    \hline
    $817$ & $431$ & $2425.31$ & $1\times10^{-4}$ & $3.208\times10^8$ \\
    \hline
    $431$ & $86$ & $2167.69$ & $1\times10^{-4}$ & $2.868\times10^8$ \\
    \hline
    \end{tabular}
    \caption{Tabella che contiene i dati effettuati facendo ruotare lo specchio in senso antiorario.}
    \label{antiorario}
\end{table}
Svolgendo la media tra le misure si ottiene:
\begin{equation}
    \nonumber
    c_{spe}=2.997\times10^8 \, \frac{m}{s} \qquad \sigma_c=7.8\times10^6 \, \frac{m}{s} 
\end{equation}
Per valutare quantitativamente la misura ottenuta rispetto al valore teorico di $c$ che nell'aria equivale a circa: $c_{teo}=2.998\times10^{8} \,ms^{-1}$ è stato effettuato il test t di Student con $N-1$ gradi di libertà dove $N$ è il numero che corrisponde alle misure effettuate:
\begin{equation}
    t=\frac{|c_{spe}-c_{teo}|}{\sigma_c}=0.015
\end{equation}
Dopo aver consultato le opportune tabelle si ottiene che c'è $98\%$ di probabilità che la misura sia diversa dal valore teorico per fluttuazioni statistiche.
\newpage
\subsection{Modalità turbo}
Successivamente, abbiamo preso un secondo set di dati con il "turbo". Questo metodo consiste nel sfruttare una modalità del motore dello specchio rotante che permetteva di raggiungere alte frequenze di rotazioni se pur in un breve periodo di tempo, dovuti al surriscaldamento. Di seguito sono presenti tutti i dati presi utilizzando questo secondo metodo.
\begin{table}[ht]
    \centering
    \begin{tabular}{|c|c|c|c|c|}
    \hline
     $\nu_0$ [$Hz$] & $\nu$ [$Hz$] & $\Delta \omega$ [1/s] & $\Delta\delta[m]$ & $c_{spe}$ [$m/s$] \\
    \hline
    $1409$ & $-1487$ & $18705.04$ & $8.8\times10^{-4}$ & $2.812\times10^8$ \\
    \hline
    $1492$ & $-1497$ & $18780.44$ & $8.9\times10^{-4}$ & $2.792\times10^8$ \\
    \hline
    $1496$ & $-1508$ & $18874.69$ & $8.55\times10^{-4}$ & $2.920\times10^8$ \\
    \hline
    $1480$ & $-1493$ & $18679.91$ & $7.4\times10^{-4}$ & $3.339\times10^8$ \\
    \hline
    $1500$ & $-1523$ & $18994.07$ & $8.1\times10^{-4}$ & $3.102\times10^8$ \\
    \hline
    $1512$ & $-1509$ & $18981.50$ & $8.6\times10^{-4}$ & $2.919\times10^8$ \\
    \hline
    $1513$ & $-1500$ & $18931.24$ & $7.5\times10^{-4}$ & $3.339\times10^8$ \\
    \hline
    $1525$ & $-1518$ & $19119.73$ & $8.6\times10^{-4}$ & $2.941\times10^8$ \\
    \hline
    $1523$ & $-1504$ & $19019.20$ & $8\times10^{-4}$ & $3.145\times10^8$ \\
    \hline
    $1531$ & $-1510$ & $19107.17$ & $8.4\times10^{-4}$ & $3.009\times10^8$ \\
    \hline
    $1527$ & $-1498$ & $19006.64$ & $8.3\times10^{-4}$ & $3.029\times10^8$ \\
    \hline
    $1514$ & $-1506$ & $18975.22$ & $8.9\times10^{-4}$ & $2.821\times10^8$ \\
    \hline
    $1535$ & $-1508$ & $19119.73$ & $8.5\times10^{-4}$ & $2.976\times10^8$ \\
    \hline
    $1525$ & $-1504$ & $19031.77$ & $8.5\times10^{-4}$ & $2.962\times10^8$ \\
    \hline
    $1516$ & $-1513$ & $19031.77$ & $8.4\times10^{-4}$ & $2.997\times10^8$ \\
    \hline
    $1531$ & $-1513$ & $19126.02$ & $8.4\times10^{-4}$ & $3.012\times10^8$ \\
    \hline
    $1532$ & $-1514$ & $19138.58$ & $8.5\times10^{-4}$ & $2.979\times10^8$ \\
    \hline
    $1538$ & $-1527$ & $19257.96$ & $8.4\times10^{-4}$ & $3.033\times10^8$ \\
    \hline
    \end{tabular}
    \caption{Tabella che contiene i dati effettuati facendo ruotare lo specchio in modalità turbo.}
    \label{turbo}
\end{table}\\
Attraverso l'analisi dati già illustrata si ottiene:
\begin{equation}
    \nonumber
    c_{spe}=3.007\times10^8 \, \frac{m}{s} \qquad \sigma_c=4.3\times10^6\, \frac{m}{s}
\end{equation}
Si può ora verificare la consistenza del valore ottenuto con il valore teorico attraverso l'analogo test di compatibilità t di Student ottenendo con $17$ gradi di libertà:
\begin{equation}
    t=\frac{|c_{spe}-c_{teo}|}{\sigma_c}=0.21
\end{equation}
La probabilità che la misura sperimentale sia diversa dal valore teorico per fluttuazioni statistiche è del $84\%$.
\subsection{Senso orario}
Durante la fase sperimentale abbiamo riscontrato diversi problemi legati alla strumentazione messa a disposizione. Prendendo i dati in senso orario il display presente sul macchinario che gestiva la frequenza di rotazione dello specchio, non segnalava la corretta frequenza quando si superavano i $400\, Hz$. Ci si è accorti di questo fenomeno poiché si riscontrava facilmente che ad un aumento della frequenza di rotazione (testimoniata anche da un rumore più intenso prodotto dal motorino), corrispondeva una diminuzione significativa della frequenza riportata dal display. \'E stato possibile prendere qualche dato ma essendo pochi non hanno nessuna coerenza statistica che permette lo studio o l'analisi dati. 
\subsection{Confronto tra i dati ottenuti}
\'E possibile confrontare i due valori ottenuti utilizzando nuovamente il test del t di Student. In particolare, per questo tipo di confronti si utilizza la formula
\begin{equation}
    t=\frac{|c_1-c_2|}{\sqrt{\sigma_{c_1}^2+\sigma_{c_2}^2}}
\end{equation}
dove $c_1$ e $c_2$ sono i valori ottenuti dalle due esperienze e $\sigma_{c_1}$ e $\sigma_{c_2}$ i relativi errori. Il numero dei gradi di libertà in tale contesto è pari a $gdl=N_1+N_2-2$, dove $N_1$ e $N_2$ sono i numeri di dati acquisiti da ciascuna delle due esperienze.
Si ha dunque t=0.11 con 33 gradi di libertà, che corrisponde ad una probabilità del 91\% che la differenza ottenuta tra i due valori abbia origine statistica.
